\chapter*{Preface}
\emph{Order and Chance} is a journey through three territories of mathematics, number theory, algebra and
probability, from the first principles to the questions that researchers ask today. Its destination is the place
where the three meet: the statistical behaviour of prime numbers, where the most rigid objects of mathematics turn
out to obey the laws of chance. Along the way every statement is proved. Nothing is asked to be taken on trust
that the series does not prove at some point, and the few foundations assumed here are named as such.

Mathematical maturity means learning to prove, to calculate, to solve problems,
and to think like a mathematician. This volume develops those habits through
worked arguments, experiments, and problems whose solutions explain how to find
the decisive idea.

This first volume lays the ground, in eight chapters: sums, binomial coefficients and induction; complex numbers
and roots of unity; polynomials; inequalities; sequences, series and limits; integrals and Taylor's formulas;
linear algebra; and counting, with the first steps in probability. These are the tools every later tome uses, and
they are beautiful in their own right: the roots of unity turn a sum of binomial coefficients into a geometric
picture, a single inequality between means drives a dozen optimisation problems, and a determinant measures how a
linear map stretches volume.

Each chapter opens with a story: who asked the question, what they were trying to do, and what happened next, because
the mathematics is easier to remember, and more fun, when one knows that Montmort was shuffling cards, that Cassini was
measuring the solar system, and that Marilyn vos Savant received ten thousand letters for being right. Every date and
anecdote is sourced in the reading list at the end of the chapter, and the people are shown in portraits, all in the
public domain or under a free licence, credited where they appear. Then a section called ``from scratch'' rebuilds the
chapter's tools from first principles, so that nothing is assumed from a course once taken and since forgotten: how to
write a sum, what a limit is, what a determinant measures, what a probability space is. The theorems follow, each with
its full proof, and figures drawn to scale show what the formulas mean.

Most of the pages are exercises, because mathematics is learned by doing it. They are graded \textbf{$\star$} for a
warm-up, \textbf{$\star\star$} for a good afternoon and \textbf{$\star\star\star$} for gems, and every one has a hint
and a complete solution at the end of its chapter. Each solution begins with \emph{How to find it}, which explains what
to notice and why the method is natural, before the \emph{Write-up}: the aim is to learn how a proof is found, not only
how it reads. ``Pause and think'' boxes ask you to guess before you read on; the guess is the point, not whether it is
right. Each chapter closes with a challenge box, three doors into open or deep questions, with notes on where the keys
are kept.

Every theorem, proposition and lemma is proved, including those a refresher usually quotes: the intermediate value
theorem and the mean value theorem, Stirling's formula with its constant, the Cayley--Hamilton theorem, the spectral
theorem, and Bernoulli's law of large numbers. The only facts taken without proof are the foundations: the axioms of
$\N$ and of $\R$ (in particular that every bounded set of reals has a least upper bound), the existence of the integral
of a continuous function, and the basic properties of $\exp$, $\ln$, $\cos$, $\sin$ and $\pi$. Tome~6 constructs the
real numbers and proves every fact in this list after the axioms. Results quoted only for their interest appear as
remarks, with a reference to where they are proved.

Each chapter has a companion Python notebook, in which the identities of the chapter are tested numerically and
put to work, often on a problem from computer science.

\section*{Three ways through this book}
For a gentle first journey, read each opening story, the from-scratch material and the worked examples, then attempt the one-star exercises. For a proof-building journey, add the main theorems, the new workshops and the two-star exercises. For a longer expedition, work through the three-star problems and their solutions; the challenge boxes deliberately point beyond this volume and do not all ask for elementary solutions.

Each new workshop separates the question, the hint and the complete argument. Cover the solution before trying the question; a productive wrong turn is more memorable than a proof merely watched. The field guide at the end collects the habits that connect the eight chapters.

The ordering is thematic rather than strictly axiomatic: Chapter~2 uses series justified in Chapter~5, Chapter~3 uses real completeness, and Chapter~4 uses differentiation justified in Chapters~5--6. These forward references are explicit dependencies, not extra hypotheses hidden inside a proof. Readers wanting the analytic foundations first may read the sequences and continuity sections of Chapter~5 and the differentiation toolkit of Chapter~6 before returning to Chapters~2--4. Finite-set counting and elementary logic are used throughout; the arithmetic bridge in Chapter~1 proves the divisibility facts needed before Tome~1.
